\documentclass[a4j]{jarticle}
\usepackage{ascmac}
\usepackage{comment}
\usepackage[dvipdfmx]{graphicx}
\usepackage{listings,jvlisting, jlisting}

\lstset{
  basicstyle={\ttfamily},
  identifierstyle={\small},
  commentstyle={\smallitshape},
  keywordstyle={\small\bfseries},
  ndkeywordstyle={\small},
  stringstyle={\small\ttfamily},
  frame={tb},
  breaklines=true,
  columns=[l]{fullflexible},
  numbers=left,
  xrightmargin=0zw,
  xleftmargin=3zw,
  numberstyle={\scriptsize},
  stepnumber=1,
  numbersep=1zw,
  lineskip=-0.5ex
}

\begin{document}

\vfill
\begin{screen}
  \begin{center}
    \huge\bf 実験　レポート1\\
    \vspace{5mm}
    \large\bf
    学籍番号　2022311042 \ 名前 神野友哉\\
    \today
  \end{center}
\end{screen}
\thispagestyle{empty}
%%目次
\tableofcontents

\newpage
% "%"同士に挟まれた行に書けばいいよ
\section{実験}
%

%
\subsection{目的}
%

%
\subsection{方法や条件}
%

%
\subsection{結果}
%

%
\subsection{考察}
%

%

\newpage
\section{課題\ 1}
1.3.1節で作成したロボットに新たに可動部品を追加し、レンダリングしなさい。ただし、まだ紹介されていない3ds\ Max\ の機能を\ 2\ つ以上利用すること
\subsection{考察}
%
写真 \ref{fig:02} は写真 \ref{fig:03}と同一です
%


% 下のコードで画像を二枚組にする。newpage は任意
\newpage
\begin{figure}[htbp]
  \begin{minipage}{0.5\hsize}
    \begin{center}
      \includegraphics[width=70mm]{image01.png}
    \end{center}
    \caption{image01.png}
    \label{fig:02}
  \end{minipage}
  \begin{minipage}{0.5\hsize}
    \begin{center}
      \includegraphics[width=70mm]{image01.png}
    \end{center}
    \caption{image01.png}
    \label{fig:03}
  \end{minipage}
\end{figure}

\newpage
\section{課題\ 2}
課題\ 1 \ で利用した\ 3ds\ Max\ の機能について、その特徴、役割、操作方法などを解説しなさい。
\subsection{考察}
%

%

\newpage
\section{課題\ 3}
この実験では人間が手作業で行う\ CG\ 制作を体験しているが、最近では生成系\ AI\ を利用して\ CG\ を製作することも盛んに行われている。人間が手作業で行う\ CG\ 製作と生成系\ AI\ が作成する
\ CG\ を比較して、人間でなければ作成できない部分、人間では作成できない部分について述べなさい。
\subsection{考察}
%

%

\newpage
\section{実験\ 1-1}
スプラインを作成するときにクリックしていく頂点を時計回り/反時計回りに選択してモデリングした時、概観、ポリゴン数、レンダリング時間などに
どういった差異が出るかを確認しなさい。ただし、レンダリング設定で\verb$[範囲]$ を\ 0 \
、\verb$[終端]$ を\ 99\ にしてレンダリングし、その\ 1 \ 枚当たりの平均時間を求めること。
\subsection{考察}
%

%

\newpage
\section{実験\ 1-2}
スプレインのステップ数を\ 2\ ～\ 20\ に\ 2\ ずつ変更してモデリングした時、概観、ポリゴン数、
レンダリング時間などにどういった差異が出るかを確認しなさい。
\subsection{考察}
%

%

\newpage
\section{実験\ 1-3}
レイズのセグメント数を\ 10\ ～\ 100\ に\ 10\ ずつ変更してモデリングした時、概観、ポリゴン数、
レンダリング時間などにどういった差異が出るかを確認しなさい。
\subsection{考察}
%

%

% 表の環境 "\\"は改行。"&"は要素区切り。"|c|c|c|"は要素中央そろえ、縦線で囲む。"\hline"は横線
\begin{table}
  \centering
  \begin{tabular}{|c|c|c|}
  \hline
  セグメント数 & ポリゴン数 & レンダリング時間 \\
  \hline
  なんか & なんか & なんか \\
  なんか & なんか & なんか \\
  なんか & なんか & なんか \\
  なんか & なんか & なんか \\
  なんか & なんか & なんか \\
  なんか & なんか & なんか \\
  \hline
  \end{tabular}
  \caption{実験\ 1-3}
\end{table}

\begin{comment}
\begin{lstlisting}[caption= 実験1-2,label=c-1-2]

\end{lstlisting}
\lstinputlisting[]{question4next}
\end{comment}

%%下のは参考文献. \begin{thebibliography}{3}の3は作りたい分の\bibitemの数を入れる
\newpage
\begin{thebibliography}{3}

\bibitem{lamport94}
Leslie Lamport,
\textit{LaTeX: A Document Preparation System},
Addison-Wesley, Reading, Massachusetts, 2nd Edition, 1994.

\bibitem{mittelbach04}
Frank Mittelbach, Michel Goossens, Johannes Braams, David Carlisle, and Chris Rowley,
\textit{The LaTeX Companion},
Addison-Wesley, Reading, Massachusetts, 2nd Edition, 2004.

\bibitem{knuth84}
Donald E. Knuth,
\textit{The TeXbook},
Addison-Wesley, Reading, Massachusetts, 1984.

\end{thebibliography}

\end{document}
